\subsection*{Coprime local tests and the arithmetic lattice}

\statement{B.20}{Proved}{The degree-one coprime fixed group and the direction of its lifting}
Let \(P=\Gamma_{0,12}\) be the pure discrete group, and let
\(P^{(\ell)}\) be its pro-\(\ell\) completion. Conjugation by
the specified geometric \(A_5\) defines a genuine action on it.
For \(\ell\nmid60\),
\[
 C_{P^{(\ell)}}(A_5)\longrightarrow
       (P^{(\ell),\mathrm{ab}})^{A_5}\cong\mathbb Z_\ell^2
\]
is onto. Thus the fixed group is infinite. Under the continuous
geometric outer action on \(\widehat K^{(\ell)}\), every such
fixed element has the same relative pro-\(\ell\) outer action
as some element of \(\mathcal H\), obtained by R2 for \(U=K\).
This asserts lifting to \(\mathcal H\), not lifting to
\(C_{\widehat\Gamma}(A_5)=\mathcal K\).
\begin{proof}
The integral abelianization of the pure sphere braid group is
generated by unordered-pair classes \(e_{ij}=e_{ji}\), with
relations \(\sum_{j\ne i}e_{ij}=0\). Over \(\mathbb Z_\ell\)
with \(\ell\ne2\), the vertex-relation map from the permutation
module on twelve vertices is injective: a relation in its kernel
has coefficients satisfying \(a_i+a_j=0\) for every distinct
pair, which implies \(2a_i=0\). Averaging is exact because
\(60\) is a unit. The pair action has three orbits, of sizes
\(6,30,30\), while the vertex action has one orbit.
Consequently the invariant quotient has rank \(3-1=2\).
The explicit coset enumeration in
\texttt{check\_icosahedral\_groups.py} verifies the orbit sizes.

For any pro-\(\ell\) group \(L\) acted on by the finite
\(\ell'\)-group \(A_5\), the map
\(C_L(A_5)\to C_{L/N}(A_5)\) is onto for every invariant
closed normal subgroup \(N\). For a fixed coset, its lifting
obstruction is in \(H^1(A_5,N)\), which is trivial by the
profinite Schur--Zassenhaus theorem. Apply this to the closed
commutator subgroup of \(P^{(\ell)}\).

Here is the precise R2 comparison. Set
\(E_\ell=\widehat B/O^\ell(\widehat K)\). Its kernel
\(\widehat K^{(\ell)}\) over \(A_5\) is characteristic and
centerless. The action of \(A_5\) on its abelianization is
faithful, so its centralizer in \(E_\ell\) is trivial.
The extension argument of B.5 therefore extends a centralizing
outer action of a fixed braid to an automorphism of \(E_\ell\)
inducing the identity on \(A_5\). The braid's action fixes
each peripheral class with exponent one. Its conjugators can
be taken in \(\widehat K^{(\ell)}\). The normalizer of a
peripheral subgroup in \(E_\ell\) is the procyclic extension
\(C_5\times\mathbb Z_\ell\); hence fixing its fifth power
and its finite quotient fixes the cusp generator itself.
Coprime conjugacy of the order-two and order-three free factors
likewise gives conjugators in the kernel. Put this automorphism
in the \(\ell\)-component and the identity in all other relative
components. Their quotient actions and their three conjugators
agree over \(A_5\), so they define the R2 input on
\(\widehat B/R_{\rm nil}(\widehat K)\). R2 lifts it to a
cusp-preserving multiplier-one automorphism of \(\widehat B\),
and B.6 identifies its outer class with an element of
\(\mathcal H\). The construction supplies no geometric
mapping-class representative for this lift. The proposed last
inference of P6 needs exactly that missing property; a quotient
centralizer can be much larger than the image of a centralizer.
\end{proof}

\statement{B.21}{Proved}{The sixfold cyclic cover and its multiplicity-free representation}
Let \(X\) be the smooth compactification of the cyclic cover of
the twelve-punctured icosahedral sphere which sends every positive
cusp meridian to \(1\in C_6\). Its genus is \(25\).
With \(E=\mathbb Q(\zeta_6)=\mathbb Q(\sqrt{-3})\), its
primitive-character cohomology has dimension ten over \(E\).
The \(A_5\)-representation on it is
\[
 V=V_1\oplus V_4\oplus V_5,
 \qquad\dim_E V_i=i,
\]
where all three constituents are absolutely irreducible and
occur once. Its polarization has signature \((1,9)\), with
the one-dimensional constituent of one sign and the other two
of the opposite sign. The unramified cover is the fibre product
defined by \(K\cap B'\leq B\). This is a congruence
subgroup of exact level thirty, and
\(B/(K\cap B')\cong A_5\times C_6\).
\begin{proof}
The meridian assignment is well defined because twelve is
divisible by six. Riemann--Hurwitz gives
\(2g-2=6(-2)+12(6-1)=48\).
In a coordinate avoiding a branch point at infinity the cover
has equation \(y^6=f_{12}(t)\). The eigendifferentials are
\[
 t^j\,dt/y^k,\qquad 1\leq k\leq5,
                 \quad 0\leq j\leq2k-2.
\]
Their dimensions are \(1,3,5,7,9\); regularity at the finite
branch points and at infinity proves the displayed bounds, and
their total is the genus, so they form a basis. The primitive
pair therefore has Hodge dimensions \(1,9\).

The binary icosahedral action preserves its degree-twelve Klein
form. In homogeneous coordinates the cover is \(y^6=f_{12}(X,Y)\),
with \(y\) of degree two; the central binary element \(-I\)
acts trivially on this weighted projective model. Thus the
action descends to \(A_5\). Its one-dimensional Hodge summand
is trivial, and the opposite Hodge summand is
\(\operatorname{Sym}^8(\mathbb C^2)\), restricted to the binary
group and descended to \(A_5\). Its character gives
\(V_4\oplus V_5\). This decomposition is defined over \(E\):
the representations of dimensions four and five have rational
models and Schur index one.

For a second, entirely explicit calculation,
\texttt{check\_icosahedral\_eigenspace.py} uses the Burau module
over \(\mathbb Z[t]/(t^2-t+1)\). Its rank-ten subquotient is
\[
 \ker(1,t,\ldots,t^{11})/\langle(1,\ldots,1)\rangle.
\]
Applying the words of B.8, and normalizing their scalar lifts
by \(\rho(\beta_2)\mapsto-\rho(\beta_2)\) and
\(\rho(\beta_3)\mapsto t^4\rho(\beta_3)\), gives an actual
\(A_5\)-representation with traces \(10,2,1,0,0\) on the
five conjugacy classes. Exact matrix multiplication enumerates
all sixty matrices and verifies the relations. These traces
are the sum of the characters of dimensions one, four and five;
the endomorphism centralizer has dimension three. The integral
matrices are supplied as data. This calculation also records
the scalar \(\rho(R_{12})=tI\), so it does not silently treat
the linear representation as a representation of the sphere
mapping class group: that representation is projective.

Finally \(B^{\rm ab}=C_6\), and \(\tau^5\) maps to a
generator of it. Since \(A_5\) is perfect, \(K\) surjects
onto \(B^{\rm ab}\). The restriction of this quotient to
\(K\), after replacing its generator by its inverse if
necessary, is exactly the meridian map defining the cyclic
cover. Its kernel is \(K\cap B'\), proving the fibre-product
description. To check its exact congruence level, reduction
modulo six gives
\[
 \operatorname{PSL}_2(\mathbb Z/6)
       \cong S_3\times A_4.
\]
The sign character on \(S_3\) and the quotient
\(A_4\to C_3\) compose with this reduction to the
abelianization \(B\to C_2\times C_3\): the images of
\(s,r\) generate the two factors. Thus
\(\Gamma(6)/\{\pm1\}\subseteq B'\), so the
intersection contains the level-thirty principal subgroup.
Conversely, if it contains a principal subgroup of level
\(N\), it contains \(\tau^N\). Membership in \(K\)
forces \(5\mid N\), and membership in \(B'\) forces
\(6\mid N\). Hence \(30\mid N\), proving exactness.
The image in \(A_5\times C_6\) is the full product:
both projections are onto and the perfect group \(A_5\)
has no nontrivial common quotient with \(C_6\).
Its order is \(360\), consistently giving twelve cusps
of width thirty and genus twenty-five.
\end{proof}

\statement{B.22}{Refuted}{Three ordinary linear characters do not exist on the stated full domain}
Lead 7's construction does not give three ordinary unit-valued
characters on every element of
\(\{\Phi\in\GTzero:\lambda\equiv\pm1\pmod5\}\).
Complex conjugation, which belongs to that domain, exchanges
the two primitive sixth-root eigenspaces. On a single
\(E\)-module its action is semilinear. The valid scalar
constructions must first restrict to \(\lambda\equiv1\pmod6\),
or have a semilinear target and a crossed multiplication law.
\begin{proof}
The cyclic covering subgroup is preserved by any peripheral
automorphism, but the induced automorphism of its deck group
is \(d\mapsto d^\lambda\). Thus for the induced action \(A\),
\(AdA^{-1}=d^\lambda\). If \(dv=\zeta_6v\), then
\(d(Av)=\zeta_6^{\lambda^{-1}}Av\). For \(c\),
\(\lambda=-1\), so this is the other primitive eigenspace.
It is distinct because \(\zeta_6\ne\zeta_6^{-1}\).
Consequently the asserted construction is not an \(E\)-linear
action on the named eigenspace on its whole stated domain.

After choosing semilinear trivializations, scalar coordinates
instead obey
\[
 \chi_i(\Phi\Psi)=\chi_i(\Phi)\,
       \sigma_{\lambda(\Phi)}\bigl(\chi_i(\Psi)\bigr),
 \qquad \sigma_\lambda(\zeta_6)=\zeta_6^\lambda.
\]
Furthermore an outer action on the punctured-sphere group
only determines a lift to the cyclic cover up to a deck
transformation. Such a transformation acts by a common sixth
root of unity on \(V\). Even on the linear subgroup, therefore,
three absolute scalar values require a specified lift; their
projective ratios are intrinsic. Finally the integral isotypic
projectors use \(1/60\). The asserted uniform direct product
over all primes cannot be justified by those projectors at
\(2,3,5\). These are distinct typing issues; none proves
vanishing of any corrected character on \(\mathcal K\).
\end{proof}

\statement{B.23}{Proved}{The real compact torus and its congruence centralizer}
Write \(\mathcal O_E=\mathbb Z[\zeta_6]\) and let
\(\Lambda\) be the rank-ten Eisenstein lattice of the
twelve-equal-weight Deligne--Mostow representation. Its projective
arithmetic group is the projectivization of \(\Aut(\Lambda)\),
with Hermitian signature \((1,9)\). The algebraic centralizer
of the icosahedral image in the projective unitary group is the
\(\mathbb Q\)-torus
\[
 T=\bigl(\operatorname{Res}^{,1}_{E/\mathbb Q}\mathbb G_m\bigr)^3
       /\operatorname{Res}^{,1}_{E/\mathbb Q}\mathbb G_m,
\]
the denominator embedded diagonally. It has dimension two and
is compact over \(\mathbb R\).

For \(\ell\nmid30\), put
\(U_\ell^1=\{u\in(\mathcal O_E\otimes\mathbb Z_\ell)^\times:
u\bar u=1\}\). The centralizer in the local congruence closure
is the projective image of
\[
 \left\{(u_1,u_4,u_5)\in(U_\ell^1)^3:
                 u_1u_4^4u_5^5\in\mu_3\right\}.       \tag{B.23.1}
\]
In particular it is infinite. This is a computation in the
congruence quotient of the lattice, not in \(\widehat\Gamma\).
\begin{proof}
The twelve weights are all \(1/6\); their sum is two.
The ordinary INT hypothesis in Deligne--Mostow's original
theorem is not satisfied: \((1-1/6-1/6)^{-1}=3/2\).
The equal-weight symmetry extension is essential here.
Thurston's treatment gives the corresponding dimension-nine
ball quotient, and Allcock's Theorem 5.1 identifies its linear
reflection lattice group with
\(\Aut(\Lambda)\), for
\(\Lambda=\Lambda_4^E\oplus\Lambda_4^E\oplus H\)
\citep{Thurston1998,Allcock2000}. We use the projectivization
when passing from braids to sphere mapping classes. In
particular the scalar recorded in B.21 is retained.

By B.21 and Schur's lemma, the algebraic linear commutant is
\(E^3\), acting by one scalar on each irreducible summand.
Preserving the Hermitian form imposes \(u_i\bar u_i=1\)
for each scalar. Projectivization divides by the common
scalar, giving \(T\). The signatures of the summands have
no effect on this scalar norm equation. Over \(\mathbb R\)
the result is \((S^1)^3/S^1\), hence compact.

At a prime not dividing sixty the rational isotypic
projectors are integral, so the lattice splits into the
three summand lattices. Each absolutely irreducible summand
has scalar centralizer even after extension of the local
field; an integral unitary scalar preserves its lattice
exactly when it belongs to \(U_\ell^1\). The determinant
of its action on \(V\) is \(u_1u_4^4u_5^5\).
Allcock's generators are triflections, so
\(\det\Aut(\Lambda)=\mu_3\). Its determinant-one
subgroup is the full integral special unitary group.
Strong approximation applies to the simply connected,
absolutely almost simple \(\mathbb Q\)-group
\(\operatorname{SU}(V)\): its real group is noncompact.
Consequently its integral points are dense in its product
of integral local groups. This is the strong approximation
criterion of \citet[Theorem 7.12]{PlatonovRapinchuk1994};
the precise statement and hypotheses are also displayed
and proved in the account
\citet[Theorem 2.3, Section 2.6]{Rapinchuk2012}.
The local closure of \(\Aut(\Lambda)\) is therefore exactly
the subgroup with determinant in \(\mu_3\). Intersecting
with the scalar commutant proves (B.23.1).

A projective element centralizing \(A_5\) also centralizes
its chosen linear lift: the possible scalar commutators
would define a character of the perfect group \(A_5\),
hence are trivial. Thus no extra projective centralizer
components were omitted. The kernel of the determinant
map on \((U_\ell^1)^3\) is a two-dimensional local torus,
already infinite, and its projective image still has
dimension two because its scalar kernel is finite.
\end{proof}

\statement{B.24}{Proved}{The exact lifting obstruction and the four-point control}
Let \(\rho:\widehat\Gamma\twoheadrightarrow L\) be any
specified compact quotient, put \(N=\ker\rho\), and suppose
\(A_5\) embeds in \(L\). For \(t\in C_L(A_5)\), choose
\(g\in\widehat\Gamma\) lifting it and define
\[
 a_t(h)=ghg^{-1}h^{-1}\in N\qquad(h\in A_5).
\]
Then \([a_t]\in H^1(A_5,N)\) is independent of the chosen
lift, and \(t\) lifts to \(\mathcal K\) if and only if
\([a_t]=1\). This applies to every element of (B.23.1).
At four points the analogous discrepancy can be computed
fully: in the modular quotient \(\widehat{C_2*C_3}\),
the centralizer of the geometric \(C_3\) is just \(C_3\),
whereas its congruence centralizer contains a one-dimensional
norm-one torus. Precisely the three elements of \(C_3\)
have trivial lifting obstruction in that quotient.
\begin{proof}
The calculation is the profinite version of B.3: changing
the lift by \(n\in N\) changes the cocycle by a
coboundary, and its vanishing is exactly the existence of
a centralizing lift. For the modular control, the profinite
free-product torsion-centralizer theorem gives
\(C_{\widehat{C_2*C_3}}(C_3)=C_3\).
Represent a generator by
\(R=\left(\begin{smallmatrix}0&-1\\1&-1\end{smallmatrix}\right)\).
For \(\ell\ne3\), the linear centralizer of \(R\)
in \(\operatorname{SL}_2(\mathbb Z_\ell)\) consists of
\(aI+bR\) with \(a^2-ab+b^2=1\). It is the norm-one
group of the quadratic \(\mathbb Z_\ell\)-algebra
\(\mathbb Z_\ell[X]/(X^2+X+1)\), and is infinite.
Density of \(\operatorname{SL}_2(\mathbb Z)\) in its local
congruence groups follows by elementary matrices and
reduction modulo \(\ell^n\). The obstruction criterion
therefore says exactly that only the image of \(C_3\)
lifts to the full profinite centralizer.

For the tetrahedral action on four punctures itself,
\(\Gamma_{0,[4]}=V_4\rtimes\operatorname{PSL}_2(\mathbb Z)\).
The geometric \(A_4=V_4\rtimes C_3\) has trivial
centralizer in the full completion: a centralizer element
projects to \(C_3\), centralizing \(V_4\) forces that
projection to be the identity, and then the \(C_3\)-fixed
part of \(V_4\) is zero. The congruence representation
forgets \(V_4\) and still has the infinite torus above.
Thus this is an explicit counterexample to the general
inference from an infinite quotient centralizer to an
infinite centralizer in the full completion.
\end{proof}

\statement{B.25}{Proved}{Solvable triangular controls and levels three and four}
For a finite solvable geometric group \(G\leq\Gamma_{0,[n]}\),
the independently proved normalizer comparison, together
with B.2, gives the centralizer comparison. If its punctured
quotient is a triangle orbifold, the discrete normalizer
modulo deck transformations is finite. The pure part of
its completed centralizer is therefore trivial. In
particular the tetrahedral and octahedral groups give
\(\mathcal K_3=\mathcal K_4=1\).
At six equal-weight points the octahedral representation
has a two-summand scalar commutant and an infinite compact
local torus, although only the identity can lift from that
torus to the full geometric centralizer.
\begin{proof}
The required mapping class groups and their finite-index
subgroups are good and of finite virtual cohomological type
in genus zero; the topology hypothesis for the geometric
normalizer is the finite-covering normalizer comparison
used in the proof of Boggi--Zalesskii, Theorem 5.1.
Its solvable induction proves that normalizer equality
without the final arbitrary finite-intersection step.
B.2 then proves the centralizer equality. The finite
three-point quotient group, followed by the torsion-free
pure mapping class group, gives the stated trivial pure
intersection. For the tetrahedral and octahedral actions,
the quotient signatures are \((0;2,3;\infty)\) and
\((0;2,3;\infty)\) with cusp stabilizers of orders three
and four before puncturing. Their deck groups are
\(A_4\) and \(S_4\), both solvable and centerless.
The same cusp-normalizer argument as B.7 makes every
multiplier-one centralizing geometric element pure,
so both displayed kernels vanish.

For six points use the cyclic triple cover
\(y^3=f_6(t)\), of genus four. The same differential
calculation as B.21 gives primitive Hodge dimensions
one and three. Under the octahedral action the
one-dimensional character and the three-dimensional
irreducible are distinct; twisting the common lift by
its one-dimensional character gives \(1\oplus V_3\).
Schur's lemma gives a projective norm-one torus of
dimension one, compact over \(\mathbb R\).
The six weights \(1/3\) satisfy INT, and
Deligne--Mostow, Theorem 11.4 and Proposition 12.7,
give its arithmetic ball lattice in \(PU(1,3)\)
\citep{DeligneMostow1986}. Its local closure contains
an open subgroup of the local special unitary group
by strong approximation; hence its centralizer contains
an open subgroup of this torus and is infinite.
The full completed centralizer of \(S_4\) is nevertheless
trivial by the preceding solvable argument. The
nonabelian cocycle of B.24 is therefore nontrivial
for every nonidentity point of this quotient centralizer.
The step available at four and six points is precisely
the solvable normalizer induction. It is unavailable
for \(A_5\), not supplied by a different torus dimension.
\end{proof}

\statement{B.26}{Open}{The arithmetic local calculation does not decide the global centralizer}
For twelve points we have not determined which elements
of (B.23.1) have trivial class in
\(H^1(A_5,\ker(\widehat\Gamma\to L_\ell))\).
Neither compactness at the real place nor strong
approximation computes this nonabelian cohomology set.
Thus Lead 6's infinite congruence centralizer is verified,
but it proves neither \(\mathcal K=1\) nor its negation.
The formulas at \(\ell\mid30\) need integral
representations with their nonsplit isotypic lattices;
the good-prime decomposition has not been extended to
those primes in this return.

\statement{B.27}{Open}{The remaining arithmetic-character calculations}
The normalized Jacobi-sum values on
\(G_{\mathbb Q(\zeta_{15})}\), the
values for the R2 lifts at \(U=K\) and level thirty,
and vanishing on \(\mathcal K\) remain undecided.
B.21 supplies the cover and the precise rational
representation; B.22 corrects the domains, the deck
ambiguity, and the integral projector issue before
any such calculation is meaningful. Since no vanishing
has been proved, the proposed implication to the two
Fibonacci representations has not been used.
